\documentclass[11pt,a4paper]{amsart}
\usepackage[T1]{fontenc}
\usepackage{lmodern,microtype,mathtools,amssymb}
\usepackage[margin=27mm]{geometry}
\usepackage{needspace}
\usepackage[hidelinks]{hyperref}
\hypersetup{pdftitle={Eventual exactness for local arc-connectivity},
 pdfsubject={An elementary proof using a maximum-outdegree partition and one saturated cut}}
\newtheorem{theorem}{Theorem}
\newcommand{\fl}[1]{\left\lfloor#1\right\rfloor}
\newcommand{\cl}[1]{\left\lceil#1\right\rceil}
\newcommand{\ar}{\to}
\newcommand{\dc}{\mathbin{\dot\cup}}
\newcommand{\Q}{Q_m}
\title[Eventual exactness for local arc-connectivity]{Eventual exactness for local arc-connectivity}
\author{}
\date{}
\begin{document}
\begin{abstract}
For each fixed $m\ge2$ and all sufficiently large $n$, the maximum number
of arcs in a simple digraph with no $m$ arc-disjoint paths between any
ordered pair is $\fl{(n+m-2)^2/4}$. We give an elementary proof using
short-path counts, a matching and one minimum arc cut.
\end{abstract}
\maketitle

Simple digraphs have no loops or parallel arcs; opposite arcs are allowed.
As in~\cite{thesis}, let $\lambda_D(u,v)$ be the maximum number of
arc-disjoint directed $u$--$v$ paths, put
$\lambda_D^{\max}=\max_{u\ne v}\lambda_D(u,v)$, and define
\[
 \ell_m^{\mathrm{dir}}(n)
 =\max\{|A(D)|:|V(D)|=n,\ \lambda_D^{\max}\le m-1\}.
\]
\begin{theorem}\label{thm:main}
For every fixed $m\ge2$, there is $N(m)$ such that
\[
 \ell_m^{\mathrm{dir}}(n)=\Q(n):=
 \fl{\frac{(n+m-2)^2}{4}}\qquad(n\ge N(m)).
\]
\end{theorem}
This proves the eventual form of Conjecture~3.1 in~\cite{thesis}.
No least threshold or formula at smaller orders is asserted.

Write $e(D)=|A(D)|$, let $e(U,W)$ count arcs from $U$ to $W$, and put
$e(U)=e(U,U)$; the sets may overlap. Subscripts restrict degrees and
neighbourhoods. For disjoint sets, $S\ar T$ means that all arcs from
$S$ to $T$ are present. Call $D$ \emph{feasible} if
$\lambda_D^{\max}\le m-1$. Such a digraph satisfies
\begin{equation}\label{eq:paths}
 \mathbf1_{uv\in A(D)}+|N^+(u)\cap N^-(v)|\le m-1\qquad(u\ne v),
\end{equation}
since the direct and two-step paths counted here are arc-disjoint.

\begin{proof}[Proof of Theorem~\ref{thm:main}]
For the lower bound, take disjoint sets $A,B$ with
$|B|=\cl{(n+m-2)/2}$ and $|A|=n-|B|$. Include all arcs $A\ar B$,
and in $B$ give each vertex the $m-2$ in-arcs from its cyclic
predecessors. For large $n$, this is a simple digraph with
$|B|(|A|+m-2)=\Q(n)$ arcs. For $a\in A,b\in B$, deleting $a\ar b$
and the $m-2$ internal in-arcs of $b$ separates the pair: no path can
enter another vertex of $A$. Paths starting in $B$ stay there and are
bounded by the target's internal indegree. All other distinct ordered
pairs are unreachable. Thus the construction is feasible.

For the upper bound, suppose there are arbitrarily large feasible
counterexamples $D$, with $e(D)>\Q(n)$. All implicit constants
below depend only on the fixed $m$. We may take $n$ sufficiently large
at each step.

\smallskip\noindent
\textbf{1. A complete forward pair outside a bounded set.}
Choose $v$ of maximum outdegree $\Delta$, and set
$B=N^+(v)$, $A=V(D)\setminus B$. By~\eqref{eq:paths},
\[
 d_B^-(z)\le m-2\ (z\in B),\qquad
 d_B^-(z)\le m-1\ (z\in A\setminus\{v\}).
\]
As $d_B^-(v)\le\Delta$, we obtain
\begin{equation}\label{eq:rough}
 e(B,V(D))\le(m-1)(n-1),\qquad
 e(D)\le |A|\Delta+(m-1)(n-1).
\end{equation}
Since $e(D)>n^2/4$, both parts have size $n/2+O(\sqrt n)$.
For $a\in A$, put $\alpha_a=\Delta-d^+(a)$ and
$\alpha=\sum_{a\in A}\alpha_a$. Then $0\le\alpha=O(n)$, since
$e(D)=|A|\Delta-\alpha+e(B,V(D))$.

For any $u$, put $W=N^+(u)\cap A$ and $h=|W|$.
Every arc leaving $W$ completes a two-step walk from $u$; there are at
most $m-1$ to each other vertex and at most $h$ back to $u$. Hence
\[
 h\Delta-\alpha\le e(W,V(D))\le(m-1)(n-1)+h,
 \qquad\text{so }d_A^+(u)=h=O(1).
\]
The missing forward targets $M(a)=B\setminus N^+(a)$ satisfy
\begin{equation}\label{eq:holes}
 |M(a)|=\alpha_a+d_A^+(a),\qquad
 \sum_{a\in A}|M(a)|=\alpha+e(A)=O(n).
\end{equation}
Averaging over $A\setminus\{v\}$ gives $u$ with $|M(u)|=O(1)$.
Counting $u\ar b\ar v$ then gives $d_B^-(v)\le m-1+|M(u)|=O(1)$.
Thus the bipartite graph of backward arcs from $B$ to $A$ has bounded degrees.

Call a vertex \emph{good} if it is incident with at most $n/(10m)$
missing forward arcs. Only $O(1)$ vertices are not good, by
\eqref{eq:holes}. A backward matching $b_i\ar a_i$ of size $m-1$
between good vertices would give $m$ arc-disjoint paths
\[
 s\ar t,\qquad s\ar b_i\ar a_i\ar t\quad(1\le i\le m-1):
\]
choose a good $s\in A\setminus\{a_i\}$ pointing to every $b_i$, then
$t\in B\setminus\{b_i\}$ receiving arcs from $s$ and all $a_i$.
Each choice excludes fewer than $n/10+O(1)$ vertices, whereas both parts
have size $n/2+O(\sqrt n)$. Such a matching is therefore impossible.
A maximal matching on good vertices has at most $m-2$ edges, and its
endpoints meet every good backward arc. Bounded degrees and the $O(1)$
nongood vertices therefore give $f:=e(B,A)=O(1)$.

Put $\beta=\sum_{b\in B}(m-2-d_B^-(b))$. The exact count
\begin{equation}\label{eq:deficit}
 e(D)=\Delta(n+m-2-\Delta)+f-\alpha-\beta
\end{equation}
implies $\alpha+\beta<f=O(1)$, since the quadratic term is at most
$\Q(n)$. In particular, every $M(a)$ is bounded.
If $\alpha_a=0$, put $p=d_A^+(a)=|M(a)|$. For all but $O(1)$
vertices $b\in B$, we have $a\ar b$, $N_A^+(a)\ar\{b\}$ and
$d_B^-(b)=m-2$. Equation~\eqref{eq:paths} yields
\[
 1+p+(m-2)-|M(a)\cap N_B^-(b)|\le m-1,
 \qquad\text{hence }M(a)\subseteq N_B^-(b).
\]
Thus every vertex of $H:=\bigcup_{\alpha_a=0}M(a)$ points to all but
$O(1)$ vertices of $B$, with a uniform bound. As $e(B)\le(m-2)|B|$,
there are at most $m-2$ vertices in $H$ for large $n$.
Take $S=\{a\in A:\alpha_a=0\}$, and obtain $T$ from $B$ by removing
$H$ and every tail of a backward arc. With $C=V(D)\setminus(S\cup T)$,
we have
\begin{equation}\label{eq:partition}
 S\ar T,\qquad e(T,S)=0,\qquad |C|=O(1),\qquad
 d^+(s)=\Delta\quad(s\in S).
\end{equation}
Both parts still have size $n/2+O(\sqrt n)$. Crucially, all vertices
of $S$ retain the same outdegree in the original digraph.

\smallskip\noindent
\textbf{2. Equal degrees force every cross-pair to be saturated.}
Write $a=|S|$, $b=|T|$, $c=|C|$, and $x=\Delta-b$.
Each core vertex has at most $m-1$ successors in $S$ and at most $m-1$
predecessors in $T$, by counting two-step paths to $T$ and from $S$.
Consequently
\[
 E:=e(C)+e(C,S)+e(T,C)\le c(c-1)+2(m-1)c=O(1).
\]
For $s\in S,t\in T$, set $I_s=N_C^+(s)$, $O_t=N_C^-(t)$ and
$\eta_t=d^-(t)-a$. Counting the direct and two-step $s$--$t$ paths gives
\begin{equation}\label{eq:effective}
 x+\eta_t
 \le m-2+|I_s\cup O_t|\le c+m-2.
\end{equation}
Let $y=\max_{t\in T}\eta_t$ and
$\delta=\sum_{t\in T}(y-\eta_t)$. Then
\begin{equation}\label{eq:count}
 e(D)=ab+ax+by+E-\delta
     =(a+y)(b+x)+E-xy-\delta.
\end{equation}
If $x+y<c+m-2$, the product is at most
$\fl{(n+m-3)^2/4}$, contradicting $e(D)>\Q(n)$ for large $n$ since
$E=O(1)$. Thus $x+y=c+m-2$ and $\delta<E-xy\le E$.
Move the $O(1)$ vertices with $\eta_t<y$ from $T$ into $C$.
If $h$ vertices move, then $x$ and $c$ both increase by $h$, while
$y$ is unchanged. Retain the notation for this new partition: $E$
remains bounded, $\delta=0$, and
\begin{equation}\label{eq:uniform}
 d^+(s)=b+x\ (s\in S),\qquad d^-(t)=a+y\ (t\in T),\qquad
 x+y=c+m-2.
\end{equation}
Equality throughout~\eqref{eq:effective} shows that, for every $s,t$,
\begin{equation}\label{eq:saturated}
 I_s\cup O_t=C,\qquad
 d_S^+(s)+d_T^-(t)+|I_s\cap O_t|=m-2.
\end{equation}
Thus their direct and two-step paths form an arc-disjoint family of
size exactly $m-1$.

\par\Needspace{8\baselineskip}
\smallskip\noindent
\textbf{3. Counting the core with one saturated cut.}
By~\eqref{eq:saturated}, $D[S]$ has maximum outdegree at most $m-2$
and $D[T]$ has maximum indegree at most $m-2$.
Choose $s\in S$ not reachable from $C$ in one step or in two steps
through $S$, and $t\in T$ not reaching $C$ in one step or in two steps
through $T$. Each condition excludes at most $(m-1)^2c$ vertices.
Put $I=I_s$, $O=O_t$, $p=d_S^+(s)$ and $q=d_T^-(t)$.

By directed max-flow/min-cut~\cite{flow}, a minimum $s$--$t$ cut has
size $m-1$. The displayed paths already require that many crossing arcs,
so \emph{every crossing arc lies on one of them}, and in particular
has tail $s$ or head $t$. For its source side $Z$,
write
\[
 P=Z\cap S,\quad R=T\setminus Z,\quad
 X=Z\cap C,\quad Y=C\setminus X.
\]
Completeness of $S\ar T$ and the cut property give
\[
 P\setminus\{s\}\subseteq N_S^+(s),\qquad
 R\setminus\{t\}\subseteq N_T^-(t),\qquad
 X\subseteq I,\quad Y\subseteq O,\quad e(X,Y)=0.
\]
For example, a vertex of $X\setminus I$ belongs to $O$ and would send
an extra crossing arc to $t$; the other core inclusion is symmetric.

For $u\in P\setminus\{s\}$, no arc leads to $(S\setminus P)\cup Y$,
since it would be an additional crossing arc. Hence
\[
 x=d^+(u)-b\le |P|-1+|X|\le p+|I|=x.
\]
Equality forces $u\ar s$. An arc from $X$ to $u$ would give a
two-step path from $C$ to $s$, contrary to its choice. That choice also
excludes arcs from $X$ to $s$, and the cut excludes arcs from $X$ to
$S\setminus P$. Dually, for $w\in R\setminus\{t\}$,
$y=d^-(w)-a\le|R|-1+|Y|\le q+|O|=y$ forces $t\ar w$;
the choice of $t$ and the cut similarly exclude all arcs from $T$ to $Y$.
Thus
\begin{equation}\label{eq:zeros}
 e(X,Y)=e(X,S)=e(T,Y)=0.
\end{equation}

For each vertex of $X\subseteq I$, count the direct path from $s$ and
the two-step paths through $I$ and $T$. Also, every internal successor
$s'$ of $s$ points to $C\setminus O$, by $I_{s'}\cup O=C$.
Summing~\eqref{eq:paths}, and then arguing dually from $Y$ to $t$, gives
\[
 \begin{aligned}
 e(I,X)+e(T,X)&\le(m-2)|X|-p(c-|O|),\\
 e(Y,O)+e(Y,S)&\le(m-2)|Y|-q(c-|I|).
 \end{aligned}
\]
By~\eqref{eq:zeros}, these sums cover all arcs counted by $E$, except
possibly those from $C\setminus I$ to $C\setminus O$; any
double counting is harmless. Since $p+|I|=x$, $q+|O|=y$ and $x+y=c+m-2$, it follows that
\[
 \begin{aligned}
 E&\le(m-2)c-p(c-|O|)-q(c-|I|)+(c-|I|)(c-|O|)\\
  &=p|O|+q|I|+|I||O|=xy-pq.
 \end{aligned}
\]
Finally~\eqref{eq:count}, now with $\delta=0$, yields
\[
 e(D)\le(a+y)(b+x)-pq\le\fl{\frac{(n+m-2)^2}{4}},
\]
because the two integer factors sum to $n+m-2$ and $p,q\ge0$.
This contradiction proves the upper bound.
\end{proof}

\begin{thebibliography}{9}
\bibitem{flow}
L.~R.~Ford, Jr. and D.~R.~Fulkerson,
\emph{Maximal flow through a network},
Canadian J. Math. \textbf{8} (1956), 399--404.
\bibitem{thesis}
L.~Vandenbruwaene,
\emph{Maximising Edge Density Subject to Connectivity Constraints:
Erd\H{o}s Problem 915, from Proof to Search},
master's thesis manuscript, KU Leuven, 2026--2027,
Conjecture~3.1 and Construction~A.17.
\end{thebibliography}
\end{document}
